\documentclass{article}
\usepackage[T1]{fontenc}
\usepackage{lmodern}
\usepackage{graphicx}
\usepackage{amsmath,amssymb}
\usepackage[a4paper,margin=2cm]{geometry}
\usepackage[absolute,overlay]{textpos}
\setlength{\TPHorizModule}{1mm}
\setlength{\TPVertModule}{1mm}
\usepackage[indonesian]{babel}
\usepackage{hyperref}
\setlength{\emergencystretch}{2em}
% unit: Halaman 11

\begin{document}

% === Halaman 11 (llm) ===

Perbedaan GOST dengan DES:

\begin{enumerate}
    \item Kunci $DES$ 56 bit, sedangkan kunci GOST lebih panjang yaitu 256 bit. Ini menyebabkan \textit{exhaustive key search} terhadap $GOST$ lebih sukar dibandingkan dengan $DES$.
    \item Jumlah putaran DES 16 kali, sedangkan GOST lebih banyak yaitu 32 kali sehingga membuat kriptanalisis menjadi sangat sulit
    \item Kotak S di dalam $DES$ menerima masukan 6 bit dan luaran 4 bit (berukuran $6 \times 4$), sedangkan kotak S di dalam GOST menerima masukan 4 bit dan luaran 4 bit (berukuran $4 \times 4$)
    \item Pembangkitan kunci internal $DES$ rumit, sedangkan di dalam $GOST$ pembangkitan kunci internalnya sederhana
    \item DES mempunyai permutasi yang tidak teratur, sedangkan GOST hanya menggunakan pergeseran 11-bit secara sirkuler
\end{enumerate}

\end{document}
